\documentclass[10pt,a4paper]{amsart}
\usepackage[margin=19mm]{geometry}
\usepackage{amsmath,amssymb,mathtools}
\usepackage[scr=rsfs,cal=euler]{mathalfa}
\usepackage{xfrac}
\usepackage[unicode,hidelinks,bookmarks=false]{hyperref}
\newcommand{\ie}{i.e.\ }
\newcommand{\cf}{cf.\ }
\newcommand{\resp}{respectively\ }
\newcommand{\Subsection}{\subsection}
\providecommand{\nobreakdash}{\nobreak\hbox{-}}
\setlength{\emergencystretch}{2em}
\multlinegap0pt
\usepackage{xcolor}
\usepackage{bm}
\usepackage{algorithm}\usepackage{algpseudocode}
\usepackage{float}
\newtheorem{lemma}{Lemma}
\newtheorem{theorem}{Theorem}
\newcommand{\E}{\mathbb{E}}
\def\gF{{\mathcal{F}}}
\def\gN{{\mathcal{N}}}
\title{The Phase Transition in Online PCA Depends on $n/d\log(d)$, not $n/d$}
\author{Apratim Dey}
\date{Source: arXiv:2607.23914v2 (CC BY 4.0)}
\begin{document}
\maketitle

\noindent\emph{Source and licence.} The Phase Transition in Online PCA Depends on $n/d\log(d)$, not $n/d$. Version arXiv:2607.23914v2 (CC BY 4.0), distributed by arXiv under the Creative Commons Attribution 4.0 International licence, http://creativecommons.org/licenses/by/4.0/, as recorded in the arXiv licence metadata for this exact version and shown on the abstract page https://arxiv.org/abs/2607.23914v2. Full licence notice: \texttt{licenses/arxiv-2607.23914-CC-BY-4.0.txt}.\par\smallskip

\noindent\textit{Editorial scope.} Counted unit: the article's lemma lemma:exact\_behavior\_rho\_subcritical (source0.tex lines 1718--1723). The article's complete original proof of it is delivered with it (source0.tex lines 1727--1786, one \texttt{proof} environment of 60 lines, which the article places away from the statement). No mathematics is written, repaired or re-derived: the statement and the proof are the article's own lines, in the article's own notation, and no retained line is substituted or re-flowed.  Retained uncounted as the source-local context the proof needs: lines 741--746; lines 790--795; lines 1561--1566; lines 1621--1626; lines 1658--1660.  Nothing was omitted from any retained span, so the package carries no omission record.  No retained line contains a citation, so the wrapper carries no bibliography and no bibliography entry.  No retained line contains a figure environment, an image-inclusion command or drawing code, so no image is delivered and the package declares no asset dependency.  Editorial formatting only, in addition: an AMS \texttt{amsart} preamble that restates the article's own declarations and macros used by the retained lines, namely the environments \texttt{lemma}, \texttt{theorem} on separate counters, together with the standard packages those lines need.  The retained environments print in this document as Lemma 1, Theorem 1 in order of appearance, whereas the article prints the same items with the numbering of its own full document.  The complete native source of the article as distributed in the arXiv e-print for this exact version is bundled unchanged as \texttt{sources/206075/source0.tex}.
\begin{lemma}\label{lemma:Xk_components_indep}
The random variables $A_k, C_k, \|D_k\|^2$ are mutually independent, and also independent of $\gF_{k-1}$. Further,
\begin{align*}
    A_k\sim \gN(0,\theta^2+1), \quad C_k\sim \gN(0,1),\quad \|D_k\|^2\sim \chi^2_{d-2}
\end{align*}
\end{lemma}

\begin{theorem}\label{thm:subcritical}(Subcritical phase)
    Suppose $n/d\log d\to \gamma$ as $d\to\infty$. Then,
    \begin{align*}
        \gamma<\gamma_*\implies \rho_n\stackrel{p}{\to}0
    \end{align*}
\end{theorem}

\begin{align}\label{eqn:martingale_def}
\begin{split}
    M_k &= \delta\rho_{k-1}(A_k^2-(\theta^2+1))+\delta\sqrt{1-\rho^2_{k-1}}A_kC_k-\delta\rho_{k-1}(B_k^2-(\theta^2\rho_{k-1}^2+1))\\
    &-\delta^2\rho_{k-1}(B_k^2\|D_k\|^2/d-(\theta^2\rho_{k-1}^2+1))/2
\end{split}
\end{align}Clearly $M_k$ is $\gF_k$-measurable and one can check that $\E[M_k|\gF_{k-1}]=0$, using Lemma \ref{lemma:Xk_components_indep}. Thus $\{M_k\}_{k\geq 1}$ a martingale difference sequence. Further, from the expression (\ref{eqn:martingale_def}), it is clear that $M_k$ is a polynomial in the random variables $A_k,B_k,C_k,D_k^2/d$ which all have finite moments of all orders. This implies that each moment $\E|M_k|^j\leq C_j$ for all $k$. Some algebra yields the final expression, and the proof is complete. The unfolded result is a simple consequence.

\begin{lemma}\label{lemma:martingale_simplified}
    Recall the martingale difference $M_k$ defined in (\ref{eqn:martingale_def}). Let $\tilde M_k=\delta A_kC_k$. Then, if $\limsup_{d\to\infty} n/d\log d\leq \gamma_*$,
    \begin{align*}
        \dfrac{1}{d}\sum_{k=1}^n c_d^{n-k}(M_k-\tilde M_k)\stackrel{p}{\to}0
    \end{align*}
\end{lemma}

\begin{align}\label{eqn:recursion_unfolded_final_simplified}
    \rho_k &= c_d^k\rho_0 - \dfrac{\beta}{d}\sum_{k=1}^n c_d^{n-k}\rho^3_{k-1} + \dfrac{1}{d}\sum_{k=1}^n c_d^{n-k}\tilde M_k + \dfrac{1}{d}\sum_{k=1}^n c_d^{n-k}S_k + \dfrac{1}{d^2}\sum_{k=1}^n c_d^{n-k}R_k
\end{align}

\begin{lemma}\label{lemma:exact_behavior_rho_subcritical}
    Let $n=[\gamma_*d\log(d)-t_dd]$ where $t_d\to\infty,t_d/\log(d)\to0$ as $d\to\infty$. Then,
    \begin{align*}
        \exp(\alpha t_d)\rho_n\stackrel{d}{\to}\gN(0, 1+\gamma_*\delta^2(\theta^2+1))
    \end{align*}
\end{lemma}

\begin{proof}[Proof of Lemma \ref{lemma:exact_behavior_rho_subcritical}]
    We start with recursion \ref{eqn:recursion_unfolded_final_simplified}, and multiply both sides by $\exp(\alpha t_d)$. Then, following the proof of Theorem \ref{thm:subcritical},
    \begin{align*}
        \E\left|\dfrac{\exp(\alpha t_d)}{d}\sum_{k=1}^n c_d^{n-k}\rho^3_{k-1}\right|&\leq \dfrac{C\exp(\alpha t_d+3\alpha n/d)}{d^{3/2}}\\
        &\leq \dfrac{C\exp(\alpha t_d + 3\log(d)/2-3\alpha t_d)}{d^{3/2}}\\
        &=C\exp(-2\alpha t_d)\to 0
    \end{align*}This implies that 
    \begin{align*}
        \dfrac{\exp(\alpha t_d)}{d}\sum_{k=1}^n c_d^{n-k}\rho^3_{k-1}\stackrel{p}{\to}0
    \end{align*}
    Hence, even after scaling up by $\exp(\alpha t_d)$, the ``cubic" term involving $\rho_k^3$ (for $k\leq n-1$) is asymptotically negligible. Similarly, again from the proof of Theorem \ref{thm:subcritical},
    \begin{align*}
        \E\left|\dfrac{\exp(\alpha t_d)}{d^2}\sum_{k=1}^n c_d^{n-k}R_k\right| &\leq \dfrac{C\exp(\alpha t_d+\log(d)/2-\alpha t_d)}{d}= \dfrac{C}{\sqrt{d}}\to 0
    \end{align*}and thus $$\exp(\alpha t_d)\sum_{k=1}^n c_d^{n-k}R_k/d^2\stackrel{p}{\to}0$$
    Finally, we show that
    \begin{align*}
        \dfrac{\exp(\alpha t_d)}{d}\sum_{k=1}^n c_d^{n-k}S_k\stackrel{p}{\to}0
    \end{align*}To see this, we follow the proof of Lemma \ref{lemma:martingale_simplified}, and obtain
    \begin{align*}
        \E\left(\dfrac{\exp(\alpha t_d)}{d}\sum_{k=1}^n c_d^{n-k}S_k\right)^2 &= \dfrac{\exp(2\alpha t_d)}{d^2}\sum_{k=1}^n c_d^{2(n-k)}\E(S_k^2)\\
        &\leq \dfrac{C\exp(2\alpha t_d)}{d^3}\sum_{k=1}^n \exp(2\alpha(n-k)/d)\times \exp(2\alpha k/d)\\
        &\leq \dfrac{Cn\exp(2\alpha n/d)}{d^3}
    \end{align*}Since $n\leq \gamma_*d\log d$, it follows that the right side is further bounded by $C\log(d)/d\to 0$ and hence we get the desired result.

    Therefore, we are left with the ``initial" term $\exp(\alpha t_d)c_d^n\rho_0$ and the ``independent sum" term $\exp(\alpha t_d)\sum_{k=1}^n c_d^{n-k}M_k/d$. We will show that they converge to independent $N(0,1)$ and $N(0, \gamma_*(\theta^2+1)\delta^2)$. Together with the smallness of the two other terms, an application of Slutsky's lemma then implies the desired result.

    Note that
    \begin{align*}
    n\log(1+\alpha/d)+\alpha t_d-\log(d)/{2} &= n(\alpha/d + O(1/d^2)) +\alpha t_d -\log(d)/2)\\
    &=O(n/d^2)=O(\log d/d)\to 0
    \end{align*}This implies that 
    \begin{align*}
        \dfrac{\exp(\alpha t_d)\times c_d^n}{d^{1/2}}= \exp\left(n\log(1+\alpha/d)+\alpha t_d-\log(d)/{2}\right)\to 1
    \end{align*}Further, we already know that $d^{1/2}\rho_0\stackrel{d}{\to}\gN(0,1)$. Thus, we get $\exp(\alpha t_d)c_d^n\rho_0\stackrel{d}{\to}\gN(0,1)$.

    Finally, we apply the Lyapounov central limit theorem to show that
\begin{align*}
    \dfrac{\exp(\alpha t_d)}{d}\sum_{k=1}^n c_d^{n-k}\tilde M_k\stackrel{d}{\to}\gN(0,\gamma_*(\theta^2+1)\delta^2)
\end{align*} Observe that the terms $\tilde M_k=\delta A_kC_k$ are iid with mean $0$, variance $\delta^2(\theta^2+1)$ and bounded moments of all orders. Consequently, the variance of this sum yields
\begin{align*}
    \E\left(\dfrac{\exp(\alpha t_d)}{d}\sum_{k=1}^nc_d^{n-k}\tilde M_k\right)^2 &= \dfrac{\exp(2\alpha t_d)\delta^2(\theta^2+1)}{d^2}\sum_{k=0}^n c_d^{2k}\\
    &=\dfrac{\delta^2(\theta^2+1)\exp(2\alpha t_d)(\exp(2\alpha n/d)-1)}{2\alpha}
\end{align*}Now, using that $\gamma_*d\log d-t_d d-1\leq n\leq \gamma_*d\log d-t_d d$, it follows that
\begin{align*}
    \log d - 2\alpha t_d - 2\alpha/d \leq 2\alpha n/d\leq \log d - 2\alpha t_d
\end{align*}which finally implies that
\begin{align*}
    \dfrac{\exp(2\alpha t_d + 2\alpha n/d) }{d}\to 1
\end{align*}and hence, 
\begin{align*}
    \E\left(\dfrac{\exp(\alpha t_d)}{d}\sum_{k=1}^nc_d^{n-k}\tilde M_k\right)^2 \to \gamma_*\delta^2(\theta^2+1)
\end{align*} The last step is to verify the Lyapounov condition is satisfied, that is, we need to show that
\begin{align*}
    \dfrac{\exp(3\alpha t_d)}{d^3}\sum_{k=1}^nc_d^{3(n-k)}\E|\tilde M_k|^3\to 0
\end{align*}But each $\E|\tilde M_k|^3\leq C$, so that
\begin{align*}
    \dfrac{\exp(3\alpha t_d)}{d^3}\sum_{k=1}^nc_d^{3(n-k)}\E|\tilde M_k|^3 &\leq \dfrac{C\exp(3\alpha t_d)\times \exp(3\alpha n/d)}{d^2}\\
    &\leq \dfrac{C\exp(3\alpha \gamma_*\log d)}{d^2}\\
    &=\dfrac{Cd^{3/2}}{d^2}\to 0
\end{align*}and the proof is complete.\end{proof}

\end{document}
